\section*{PRÁCTICA No. 9}

\subsection*{RECUPERACIÓN DE FUENTES USO DE MANIPULADORES A DISTANCIA}

\textbf{OBJETIVO:} Confirmar, si al final del curso recibido, el alumno cuenta ya con los conocimientos suficientes como para saber manejar satisfactoriamente una situación de pérdida de control de una fuente radiactiva.

\textbf{MATERIAL:}
\begin{itemize}
    \item Detector de Radiaciones Tipo Geiger-Müller.
    \item Equipo Portátil Detector de Radiaciones con Alarma Sonora.
    \item Contenedor de Gammagrafía (Con todos sus Accesorios).
    \item Fuente simulada.
    \item Juego de Material de Acordonamiento y contienen:
    \begin{itemize}
        \item Pinzas Largas.
        \item Contenedor de Rescate.
        \item Embudo.
    \end{itemize}
    \item Madrina de Plomo.
\end{itemize}

\textbf{DESARROLLO:} Se supone que ha habido un corte o ruptura del chicote o telemático por la caída de alguna estructura metálica pesada con filos agudos, en la zona de trabajo.
El procedimiento a seguir en este caso es el siguiente:
Localizar y blindar la fuente radiactiva con la madrina, que siempre debe tenerse a la mano.
Aplicar el acordonamiento y la señalización hasta donde se tengan lecturas de rapidez de exposición menores a 0.25 mR/h, intensificando la vigilancia del tráfico de público en las zonas aledañas.